% =============================================================================
% Abstract and index terms
% =============================================================================
\begin{abstract}
Production AI agents frequently degrade after deployment, yet structured
remediation culture is largely absent. Empirical studies confirm most
deployments fail or underdeliver (Pan et al., 2025), and simulation-based
research projects that unchecked behavioral drift reduces task success by
42\% (Rath, 2026). We introduce the Agent Performance Improvement Plan
(APIP): a lifecycle protocol adapted from organizational management that
governs the remediation of underperforming agents. Drawing on human PIP
structure and Kirkpatrick's four-level evaluation framework, the APIP defines
statistically disciplined trigger thresholds, a five-category failure mode
taxonomy, a measurable cause-attribution pipeline, a tiered intervention
hierarchy, a bounded remediation window with check-ins, and exit criteria. We
evaluate the protocol in three live multi-agent environments spanning customer
service ($\tau^2$-bench), software engineering (ChatDev), and economic
operations (CoffeeBench), injecting each failure mode with known ground truth
and scoring the full pipeline against no-action, fixed-tier, random-tier, and
oracle baselines. APIP-governed remediation achieves \ttrOraclePct{} of oracle
time-to-resolution versus \ttrRandomPct{} for random-tier selection, and the
attribution pipeline identifies the injected failure mode with accuracy
\attrAccuracy{} while abstaining honestly in \attrAbstention{} of cases---to
our knowledge the first attribution confusion matrix reported for a
remediation protocol. We formalize a declarative schema meeting AI governance
traceability standards (Kolt et al., FAccT 2024; EU AI Act), and demonstrate
the mesh disruption failure mode---where a locally optimal new agent degrades
peer performance through context disruption (Kim et al., 2025)---together with
a scoped-context onboarding protocol that attenuates it, grounded in Event
System Theory and team coordination research.
\end{abstract}

\begin{IEEEkeywords}
AI agents, post-deployment monitoring, remediation, behavioral drift,
AI governance, multi-agent systems, performance improvement plan
\end{IEEEkeywords}
